\chapter{Hypothesis 2: Fluid-Front Tracking and Material-Change Decoupling}
\label{ch:hyp2}

\Cref{ch:hyp1} used an idealised target: a point-like PEC cylinder, chosen
because its response is easy to interpret and ground-truth. This chapter
applies the same migration and phase-plane machinery to a more
field-realistic scenario --- a spatially \emph{distributed} fluid front
advancing laterally through a sub-wavelength fracture, following the
thin-layer reflectivity model of \cref{sec:th-material-change}. Unlike a
point scatterer, a fluid front is not expected to move as a rigid body, so
this chapter is also the first practical test of the geometry/material-change
decoupling derived in \cref{sec:th-material-change}.

\paragraph{Hypothesis 2.} Multi-dimensional phase-plane regression allows
tracking the subwavelength translation of a fluid front and inferring the
associated material change in a subwavelength thin fracture.

\section{Setup: Raw B-Scans and Signal Conditioning}

The fluid-front model reuses the domain, grid, and Ricker source of
\cref{ch:methodology}; the fracture is represented as a thin, sub-wavelength
layer whose air-filled and water-filled portions are separated by a front
that advances laterally between the baseline and each of seven monitor
scenarios, from $2\lambda$ down to $\tfrac{1}{32}\lambda$.
\Cref{fig:ff-bscans} shows the raw and background-subtracted B-scans, and
\cref{fig:ff-taper} the effect of the standard tapering and $t_0$-shift
conditioning (\cref{sec:meth-conditioning}) on the $2\lambda$ scenario.

\begin{figure}[htbp]
  \centering
  \begin{subfigure}[t]{0.48\textwidth}
    \includegraphics[width=\linewidth]{FF_001_GPR_B-Scans__Background_Baseline_and_TimeLapsed_Models.png}
    \caption{Background \& time-lapsed}
  \end{subfigure}
  \hfill
  \begin{subfigure}[t]{0.48\textwidth}
    \includegraphics[width=\linewidth]{FF_002_GPR_B-Scans__Background_Subtracted.png}
    \caption{Background-subtracted}
  \end{subfigure}
  \caption{Raw B-scans for the fluid-front study: (a) background and
    time-lapsed models; (b) background-subtracted.}
  \label{fig:ff-bscans}
\end{figure}

\begin{figure}[htbp]
  \centering
  \begin{subfigure}[t]{0.48\textwidth}
    \includegraphics[width=\linewidth]{FF_003_Effect_of_Tapering_and_t0_Shift__2λ_dataset_single_trace.png}
    \caption{Single trace}
  \end{subfigure}
  \hfill
  \begin{subfigure}[t]{0.48\textwidth}
    \includegraphics[width=\linewidth]{FF_004_B-scan_effect_of_tapering_and_t0_shift__2λ_dataset.png}
    \caption{Full B-scan}
  \end{subfigure}
  \caption{Effect of tapering and the $t_0$ shift on the $2\lambda$
    fluid-front dataset: (a) a single trace; (b) the full B-scan.}
  \label{fig:ff-taper}
\end{figure}

\section{Kirchhoff Migration of the Fluid Front}

\Cref{fig:ff-kirchhoff} shows the Kirchhoff-migrated image for all eight
fluid-front scenarios and the resulting time-lapse difference.

\begin{figure}[htbp]
  \centering
  \begin{subfigure}[t]{0.48\textwidth}
    \includegraphics[width=\linewidth]{FF_005_Kirchhoff_Migration__All_8_Datasets____f_c15_GHz____aperture.png}
    \caption{All 8 scenarios}
  \end{subfigure}
  \hfill
  \begin{subfigure}[t]{0.48\textwidth}
    \includegraphics[width=\linewidth]{FF_006_Kirchhoff_Migration_zoomed____f_c15_GHz____aperture40.png}
    \caption{Zoomed}
  \end{subfigure}

  \vspace{0.6em}
  \begin{subfigure}[t]{0.48\textwidth}
    \includegraphics[width=\linewidth]{FF_007_Kirchhoff_Migration__TimeLapse_Differences_migrated__migrate.png}
    \caption{Time-lapse difference}
  \end{subfigure}
  \hfill
  \begin{subfigure}[t]{0.48\textwidth}
    \includegraphics[width=\linewidth]{FF_008_Kirchhoff_Migration__TimeLapse_Differences_zoomed.png}
    \caption{Difference, zoomed}
  \end{subfigure}
  \caption{Kirchhoff migration of the fluid-front study
    ($f_c=\SI{15}{\GHz}$, aperture $=40$): (a,b) migrated image for all
    scenarios and zoomed; (c,d) time-lapse difference and zoomed.}
  \label{fig:ff-kirchhoff}
\end{figure}

\section{Gazdag Migration of the Fluid Front}

\Cref{fig:ff-gazdag} repeats the analysis with Gazdag phase-shift migration.

\begin{figure}[htbp]
  \centering
  \begin{subfigure}[t]{0.48\textwidth}
    \includegraphics[width=\linewidth]{FF_009_Gazdag_Phase-Shift_Migration__All_7_Datasets____f_c15_GHz.png}
    \caption{All scenarios}
  \end{subfigure}
  \hfill
  \begin{subfigure}[t]{0.48\textwidth}
    \includegraphics[width=\linewidth]{FF_010_Gazdag_Phase-Shift_Migration_zoomed____f_c15_GHz.png}
    \caption{Zoomed}
  \end{subfigure}

  \vspace{0.6em}
  \begin{subfigure}[t]{0.48\textwidth}
    \includegraphics[width=\linewidth]{FF_011_Gazdag_Migration__TimeLapse_Differences_migrated__migrated_b.png}
    \caption{Time-lapse difference}
  \end{subfigure}
  \hfill
  \begin{subfigure}[t]{0.48\textwidth}
    \includegraphics[width=\linewidth]{FF_012_Gazdag_Migration__TimeLapse_Differences_zoomed.png}
    \caption{Difference, zoomed}
  \end{subfigure}
  \caption{Gazdag phase-shift migration of the fluid-front study
    ($f_c=\SI{15}{\GHz}$): (a,b) migrated image and zoomed; (c,d) time-lapse
    difference and zoomed.}
  \label{fig:ff-gazdag}
\end{figure}

\section{Back-Propagation Migration of the Fluid Front}

\Cref{fig:ff-backprop} shows the back-propagation result for both the
field-magnitude and $E_z$ images, focused at $t=\SI{1906}{\ns}$, and the
$E_z$ time-lapse difference.

\begin{figure}[htbp]
  \centering
  \begin{subfigure}[t]{0.48\textwidth}
    \includegraphics[width=\linewidth]{FF_013_Back-Propagation_E__All_7_Datasets____focus_at_1906_ns.png}
    \caption{$\lVert\mathbf{E}\rVert$, all scenarios}
  \end{subfigure}
  \hfill
  \begin{subfigure}[t]{0.48\textwidth}
    \includegraphics[width=\linewidth]{FF_014_Back-Propagation_E_zoomed____focus_at_1906_ns.png}
    \caption{$\lVert\mathbf{E}\rVert$, zoomed}
  \end{subfigure}

  \vspace{0.6em}
  \begin{subfigure}[t]{0.48\textwidth}
    \includegraphics[width=\linewidth]{FF_015_Back-Propagation_Ez__All_7_Datasets____focus_at_1906_ns.png}
    \caption{$E_z$, all scenarios}
  \end{subfigure}
  \hfill
  \begin{subfigure}[t]{0.48\textwidth}
    \includegraphics[width=\linewidth]{FF_016_Back-Propagation_Ez_zoomed____focus_at_1906_ns.png}
    \caption{$E_z$, zoomed}
  \end{subfigure}

  \vspace{0.6em}
  \begin{subfigure}[t]{0.48\textwidth}
    \includegraphics[width=\linewidth]{FF_017_Back-Propagation__TimeLapse_Differences_Ez_Ez__Ez_baseline.png}
    \caption{$E_z$ time-lapse difference}
  \end{subfigure}
  \hfill
  \begin{subfigure}[t]{0.48\textwidth}
    \includegraphics[width=\linewidth]{FF_018_Back-Propagation__TimeLapse_Differences_Ez_zoomed.png}
    \caption{Difference, zoomed}
  \end{subfigure}
  \caption{Time-reversal back-propagation migration of the fluid-front
    study, focused at $t=\SI{1906}{\ns}$: field-magnitude (a,b) and $E_z$
    (c,d) images, and the $E_z$ time-lapse difference (e,f).}
  \label{fig:ff-backprop}
\end{figure}

\section{Amplitude-Based Detectability Summary}
\label{sec:ff-detectability}

\Cref{fig:ff-summary} overlays the signed time-lapse-difference amplitude
from all three algorithms and plots the normalised lateral PSF of that
difference at the true front location, exactly as in
\cref{sec:tl-detectability} for the point-scatterer case.

\begin{figure}[htbp]
  \centering
  \begin{subfigure}[t]{0.48\textwidth}
    \includegraphics[width=\linewidth]{FF_019_TimeLapse_Migration_Comparison__Signed_Amplitude____f_c15_GH.png}
    \caption{Signed-amplitude comparison}
  \end{subfigure}
  \hfill
  \begin{subfigure}[t]{0.48\textwidth}
    \includegraphics[width=\linewidth]{FF_020_Normalised_Lateral_PSF__TimeLapse_Difference_at_True_Scatter.png}
    \caption{Normalised lateral PSF of the difference}
  \end{subfigure}
  \caption{(a) Signed time-lapse-difference amplitude for all three
    migration algorithms; (b) normalised lateral PSF of the difference
    image at the true front location.}
  \label{fig:ff-summary}
\end{figure}

\section{Phase-Plane Fit Applied to the Fluid Front}
\label{sec:ff-phaseplane}

The fluid front is not a rigid point scatterer, so the phase-plane fit here
tests something qualitatively different from \cref{ch:hyp1}: rather than
recovering a single $(\Dz,\Dx)$ displacement, the goal is to confirm that
the fit correctly attributes the front's advance to the
\emph{material-change} channel $c=\Dtheta$ of \cref{sec:th-material-change}
rather than to a spurious geometric shift, since the fracture walls
themselves do not move. \Cref{fig:ff-phaseplane} applies the estimator,
cropped to $\pm25\lambda$ around the front position located from the
smallest-shift scenario's difference-envelope peak, to all three migration
methods.

\begin{figure}[htbp]
  \centering
  \begin{subfigure}[t]{0.32\textwidth}
    \includegraphics[width=\linewidth]{TLP_013_FluidFlow__Kirchhoff____base_vs_mon_real____crop_25λ__baseli.png}
    \caption{Kirchhoff}
  \end{subfigure}
  \hfill
  \begin{subfigure}[t]{0.32\textwidth}
    \includegraphics[width=\linewidth]{TLP_014_FluidFlow__Gazdag____base_vs_mon_real____crop_25λ__baseline.png}
    \caption{Gazdag}
  \end{subfigure}
  \hfill
  \begin{subfigure}[t]{0.32\textwidth}
    \includegraphics[width=\linewidth]{TLP_015_FluidFlow__Back-prop____base_vs_mon_real____crop_25λ__baseli.png}
    \caption{Back-propagation}
  \end{subfigure}
  \caption{2D phase-plane fit applied to the fluid-front baseline/monitor
    pairs (ROI crop $\pm25\lambda$), for all three migration methods.}
  \label{fig:ff-phaseplane}
\end{figure}

\draftnote{state, for each migration method in \cref{fig:ff-phaseplane},
whether the fitted slopes $(\Dz,\Dx)$ remain near zero while the intercept
$c$ tracks the front advance, as predicted by \cref{eq:intercept-material}
--- this is the key result of the fluid-flow application and should be
stated explicitly rather than left implicit in the figure. Recall also the
centroid correction noted in the source notebook: because the newly-flooded
region spans $[x_0, x_0+\Dx_{\mathrm{true}}]$, its centroid sits at
$x_0+\Dx_{\mathrm{true}}/2$, so a raw centroid-shift readout must be doubled,
$\Dx_{\mathrm{true}} = 2\,\Dx_{\mathrm{raw}}$, to recover the true front
displacement --- confirm whether this correction is already baked into the
values plotted in \cref{fig:ff-phaseplane} or must be applied separately.
Noisy fluid-front data, where it exists, is presented in \cref{ch:hyp3}
rather than here.}
